\section{The simultaneous polynomial construction}
\label{sec:construction}

Write $U_i$ for the $i$th term of the sequence $b_0, b_1, q_0, b_2, q_1, \ldots$, beginning with $i = 1$.
At stage $i$ we shall construct $U_i = F_i / B_i$, where $F_i \in \mathbb{Q}[t]$ is monic irreducible and $B_i$ is a power of $F_1 \cdots F_{i - 1}$.
For the first stage take $B_1 = 1$.
The issue is that the subring $T$ to be studied later is arbitrary.
Lemma~\ref{lem:extraction} will produce elements of $T$ which are products of powers of the $U_h$.
The exponents are bounded in terms of one element of $T$, and the exponent of the last generator occurring is positive.
We can therefore prescribe the behaviour of all possible products in advance, using at stage $i$ the finite list
\[
    \mathcal{W}_i
    = \left\{U_i^{a} \prod_{h < i} U_h^{k_h}
      \ \middle|\ a, k_h \in \mathbb{Z},\ 1 \le a \le i,\ |k_h| \le i\right\}.
\]
Here and below a retained geometric point is a point of $\overline{\mathbb{Q}}$ at which every $F_i$ is nonzero.

The pole and separation conditions below will recover $\mathbb{Q}T$ in Section~\ref{sec:curve}.
The two arithmetic conditions have different jobs in Section~\ref{sec:conductor}.
Coefficient congruences put a power of $(t - \alpha_p)(t - \beta_p)$ in the local conductor, and large derivatives remove repeated factors.
All derivatives are formal derivatives with respect to $t$.

\begin{proposition}[Simultaneous realization]
    \label{prop:construction}
    The polynomials $F_i$, the powers $B_i$, and the arguments $\alpha_p, \beta_p$ can be chosen with all of the following properties.
    \begin{enumerate}
        \item The $U_i$ have the primary valuations in \eqref{eq:profiles} and valuation zero at every guard.
        For each nonzero $f \in \mathbb{Q}(t)$, its primary valuations eventually equal $v_E(f)$, and its guard valuations vanish with finitely many exceptions.
        The elements $a_0 = t$ and $a_{n + 1} = a_n b_{n + 1} / q_n$ belong to $D$ and satisfy $a_n - 1 \in mD$ eventually for each integer $m \ge 1$.
        \item Each $w \in \mathcal{W}_i$ has the form
        \begin{equation}
            \label{eq:scheduled-word}
            w = F_i^{a} / H_{i, w},
        \end{equation}
        where $H_{i, w}$ is a monic product of earlier polynomials, fixed before $F_i$ is chosen.
        Each $w$ has a pole at infinity and at every root of every earlier $F_h$.
        \item For each pair of distinct retained geometric points $\lambda, \mu$, all $w \in \mathcal{W}_i$ satisfy $w(\lambda) \ne w(\mu)$ once $i$ is sufficiently large.
        For each retained $\lambda$, eventually every such $w$ satisfies $w'(\lambda) \ne 0$.
        At either argument $\gamma$ over any fixed prime $p$,
        \[
            \max_{w \in \mathcal{W}_i} \nu_p(w'(\gamma)) \longrightarrow -\infty.
        \]
        \item Fix a prime $p$ and an integer $s \ge 0$.
        Set $A_s = \mathbb{Z}_p[p^s t]$, $L_p(t) = (t - \alpha_p)(t - \beta_p)$ and $g = p^{2s} L_p(t) \in A_s$.
        For every sufficiently large $i$ and every listed word $w = F_i^a/H_{i, w}$,
        \begin{equation}
            \label{eq:transfer-module}
            g^a \in p^i (A_s + A_s w).
        \end{equation}
        Thus $g^a/p^i$ is a sum $u + v w$ with $u, v \in A_s$.
        This formulation will allow us to use the estimate in any subring containing $A_s$ and $w$.
    \end{enumerate}
\end{proposition}

\begin{proof}
    \emph{The denominator and degree.}
    Suppose the preceding stages are complete.
    Express each of the finitely many products $\prod_{h < i} U_h^{k_h}$ on our list as a monomial in the preceding $F_h$ with integer exponents.
    Take $B_i = (F_1 \cdots F_{i - 1})^{K_i}$, with $K_i$ larger than the absolute values of all these exponents.
    If the older product has exponent $m_h$ on $F_h$, then $H_{i, w}$ has exponent $aK_i - m_h > 0$ on it.
    For example, at stage two we may take $K_2 = 3$.
    The listed words are $U_2^a U_1^k = F_2^a/F_1^{3a - k}$ with $1 \le a \le 2$ and $|k| \le 2$.
    Even the largest positive older exponent leaves a factor of $F_1$ in the denominator.
    Choose an integer $d_i \ge 5$ greater than all preceding degrees and large enough that $a d_i > \deg H_{i, w}$ throughout the list.
    Once we have produced an irreducible $F_i$ of degree $d_i$, it will be distinct from every predecessor, and these inequalities will give all the asserted poles.

    \emph{Assignment of the primes.}
    At the start of stage $i$, the earlier evaluation arguments are fixed.
    At every still unassigned prime the earlier polynomials have integral coefficients and unit constant terms.
    This is vacuous at stage one and will follow at each later stage from the choice of coefficient ring and constant term in \eqref{eq:raw-constant}.
    Thus every earlier polynomial is already a unit at any new primary argument in $p\mathbb{Z}_p$.
    For a new guard we must also avoid the earlier roots modulo $p$.

    Take an increasing sequence of integer cutoffs, choosing $M_i > 2^{4 + \sum_{h \le i} d_h}$ at this stage.
    Let $S_i$ consist of the primes up to $M_i$, the previously assigned primes, and the primes where the desired primary profile of $U_i$ is positive.
    For each newly assigned prime choose a transcendental element $\alpha_p \in \mathbb{Q}_p$ of the required valuation.
    Every such valuation class is uncountable, whereas the elements algebraic over $\mathbb{Q}$ are countable.
    At the first stage the finitely many primes in $S_1$ may have any nonzero guard residue.
    At each later stage, a newly assigned $p$ exceeds $M_{i - 1}$.
    Set $u = \lfloor\log_2 p\rfloor$ and choose a guard residue represented by an integer $w_p \in [u, 2u]$ at which no earlier $F_h$ vanishes modulo $p$.
    Since $p > M_{i - 1}$, this interval supplies more distinct nonzero residues than the sum of the preceding degrees.
    Choose a transcendental lift $\beta_p \in \mathbb{Z}_p$ of the selected residue.
    Both arguments are now fixed permanently.

    Choose a fresh prime $\ell > M_i$ outside $S_i$, and assign its arguments by the same rule, requiring also $w_\ell \not\equiv 1 \pmod{\ell}$.
    The cutoff leaves enough residues for this additional exclusion, including at the first stage.
    We use $\ell$ to enforce irreducibility at this stage and include it among the previously assigned primes at every later stage.
    We choose the coefficients in $\mathbb{Z}[p^{-1} : p \in S_i]$.
    The constant term $C_i = F_i(0)$ will have the form
    \begin{equation}
        \label{eq:raw-constant}
        C_i = (-1)^{d_i} \prod_{p \in S_i} p^{-N_{p, i}},
        \qquad N_{p, i} \in (\ell - 1)\mathbb{Z}_{> 0}.
    \end{equation}
    We determine the exponents $N_{p, i}$ after fixing the local targets for the two evaluations.
    Whatever these exponents are, $C_i$ is a unit outside $S_i$ and $C_i \equiv (-1)^{d_i} \pmod{\ell}$.
    At $\ell$ we impose
    \[
        F_i(t) \equiv (t - 1)^{d_i} \pmod{\ell},
        \qquad F_i(1) \equiv \ell \pmod{\ell^2}.
    \]
    The first condition is compatible with the chosen constant term.
    After the coefficients are specified modulo $\ell$, changing the linear coefficient by a multiple of $\ell$ achieves the second condition.
    These conditions make $F_i(t + 1)$ Eisenstein at $\ell$.
    They also make $F_i$ a unit at both arguments over $\ell$, since their residues are distinct from $1$.
    Its denominator $B_i$ is a unit there by our choices of arguments.

    \emph{Interpolation at the assigned arguments.}
    At each $p \in S_i$ and each $\gamma \in \{\alpha_p, \beta_p\}$, let $h$ be the required valuation of $U_i(\gamma)$, taking $h = 0$ at a guard.
    Prescribe a local target $\eta_{p, \gamma} = B_i(\gamma) p^h$ for $F_i(\gamma)$, with the following refinement when $U_i = q_n$.
    At a branch where $h = 0$ and $a_n(\gamma)b_{n + 1}(\gamma)$ is a unit, use instead
    \[
        \eta_{p, \gamma} = B_i(\gamma) a_n(\gamma) b_{n + 1}(\gamma).
    \]
    The values on the right have already been defined at that stage.
    Every target is nonzero, since evaluation at either transcendental argument is injective on $\mathbb{Q}(t)$.
    Require $F_i(\gamma)/\eta_{p, \gamma} \in 1 + p^m\mathbb{Z}_p$, taking $m = 1$ for an ordinary target and $m = \max(1, n)$ for a refined one.
    This preserves the required valuation.
    At a refined target $a_{n + 1}(\gamma) = \eta_{p, \gamma}/F_i(\gamma)$, so the same relative congruence gives $\nu_p(a_{n + 1}(\gamma) - 1) \ge n$.

    For each $p \in S_i$ let $A_p \in \mathbb{Q}_p[t]$ be the unique polynomial of degree at most one with $A_p(\gamma) = \eta_{p, \gamma} - \gamma^{d_i}$ at the two arguments.
    Then $t^{d_i} + A_p(t)$ has the prescribed values.
    Adding a scalar multiple of $L_p$ leaves both values fixed and changes the derivative at each of its two simple roots.
    To prescribe the constant term as well, use $L_p(0) \ne 0$ and put
    \begin{align*}
        J_{p, i}(t) &= t^{d_i} + A_p(t) - A_p(0)L_p(t)/L_p(0), \\
        F_{p, i}(t) &= C_i L_p(t)/L_p(0) + J_{p, i}(t).
    \end{align*}
    Here $J_{p, i}$, the targets $F_{p, i}(\gamma) = \eta_{p, \gamma}$, and all the $H_{i, w}$ are fixed independently of the exponents in \eqref{eq:raw-constant}.
    For each $p \in S_i$ choose an integer $h_p \ge 0$ such that multiplying each polynomial in the finite list
    \[
        J_{p, i}(Z/p^s), \qquad H_{i, w}(Z/p^s)
        \quad (0 \le s \le i,\ w \in \mathcal{W}_i)
    \]
    by $p^{h_p}$ gives coefficients in $\mathbb{Z}_p$.
    Take $N_{p, i} \ge h_p + i$.
    Here $N_{p, i} = -\nu_p(C_i)$, so division by $C_i$ raises each coefficient valuation by $N_{p, i}$.
    Since $N_{p, i} - h_p \ge i$, this puts both $J_{p, i}(Z/p^s)/C_i$ and every $H_{i, w}(Z/p^s)/C_i$ in $p^i\mathbb{Z}_p[Z]$.

    We may increase these exponents further before making the rational coefficient approximation.
    At either $\gamma$ over $p$ the derivative of $F_{p, i}^{a}/H_{i, w}$ is affine in $C_i$.
    The coefficient of $C_i$ in this derivative is
    \[
        \frac{a \eta_{p, \gamma}^{a - 1} L_p'(\gamma)}
             {L_p(0) H_{i, w}(\gamma)}.
    \]
    It is nonzero because the roots of $L_p$ are simple and every factor of the denominator is nonzero at $\gamma$.
    For each derivative on the finite list there is therefore a lower bound on $-\nu_p(C_i)$ beyond which the linear term has smaller valuation than the constant term and has valuation less than $-i$.
    Choose the $N_{p, i}$ to exceed all the required local bounds and to be multiples of $\ell - 1$.
    These choices are independent across $p$, since the other factors in \eqref{eq:raw-constant} are $p$-adic units.

    \emph{Rational coefficients and geometric conditions.}
    At each $p \in S_i$ require the free coefficients of $F_i$ to be sufficiently close to those of $F_{p, i}$ to preserve all the valuation, return, and derivative conditions just imposed.
    Also require for every $0 \le s \le i$ that
    \[
        \bigl(F_i(Z/p^s) - F_{p, i}(Z/p^s)\bigr) / C_i
            \in p^i \mathbb{Z}_p[Z].
    \]
    These are finitely many open conditions and contain a product of coefficient balls.
    These local conditions can be met together with the conditions at $\ell$ by coefficients in $\mathbb{Z}[p^{-1} : p \in S_i]$.
    Indeed, first multiply the local coefficient targets by a sufficiently large common product of powers of primes in $S_i$.
    The resulting conditions contain ordinary integral congruence classes at these primes and at $\ell$, so the integer Chinese remainder theorem supplies a solution after this scaling.
    Once one solution is fixed, each free coefficient may be varied independently by any integer multiple of a sufficiently divisible common integer.
    Thus the available coefficient vectors contain a translate of a full integer grid in the affine space of monic degree-$d_i$ polynomials with constant term $C_i$.
    Such a grid is Zariski dense over $\overline{\mathbb{Q}}$, by successive application of the univariate root bound.

    Enumerate the elements of $\overline{\mathbb{Q}}$ once and for all.
    Consider the first $i$ of them, discarding any root of a preceding $F_h$, and require $F_i$ to be nonzero at each remaining point.
    At the same stage require every $w$ on our list to distinguish each pair of these points and to have nonzero derivative at each of them.
    We check that each forbidden equality cuts out a proper algebraic subset of the free coefficient space.
    Nonvanishing of $F_i(\lambda)$ is automatic at zero and can be arranged at nonzero $\lambda$ by varying the linear coefficient.
    If $\lambda, \mu$ are distinct and nonzero, their evaluation rows on the coefficients of $t, t^2$ are $(\lambda, \lambda^2)$ and $(\mu, \mu^2)$, which are independent.
    Thus $F_i(\lambda)/F_i(\mu)$ varies nontrivially.
    If either point is zero, its value is the fixed nonzero constant and the other value varies.
    As the values of $H_{i, w}$ at these points are fixed and nonzero, equality of the two values of $F_i^a/H_{i, w}$ gives a proper exclusion.
    At nonzero $\lambda$ the rows $(\lambda, \lambda^2)$ and $(1, 2\lambda)$ are independent, so the logarithmic derivative $F_i'(\lambda)/F_i(\lambda)$ varies.
    At zero it is the free linear coefficient divided by $C_i$.
    Vanishing of $w'(\lambda)$ would require $aF_i'(\lambda)/F_i(\lambda) = H_{i, w}'(\lambda)/H_{i, w}(\lambda)$, again a proper exclusion.
    There are only finitely many exclusions at any one stage, and the grid meets their common complement.
    This completes the stage.
    A point on the list which is still retained at stage $i$ remains on all subsequent lists and is protected from all subsequent numerators.
    In particular every point retained in the final ring is tested at every sufficiently late stage, giving the geometric conclusions.

    \emph{The divisibility estimate.}
    Fix $p$ and $s$ as in part~4.
    At every sufficiently late stage we have $p \in S_i$ and $s \le i$.
    The formula for the interpolating polynomial gives
    \[
        \frac{F_{p, i}}{C_i} = \frac{L_p}{L_p(0)} + \frac{J_{p, i}}{C_i}.
    \]
    The coefficient bounds for $J_{p, i}$ and $H_{i, w}$ and the rational approximation to $F_{p, i}$ therefore give
    \begin{equation}
        \label{eq:local-coefficient-estimate}
        \frac{F_i(Z/p^s)}{C_i} - \frac{L_p(Z/p^s)}{L_p(0)}
        \in p^i\mathbb{Z}_p[Z],
        \qquad
        \frac{H_{i, w}(Z/p^s)}{C_i} \in p^i\mathbb{Z}_p[Z].
    \end{equation}
    Put $\lambda = p^{2s}L_p(0) \in \mathbb{Z}_p$ and substitute $Z = p^s t$.
    We have $g = (p^s t - p^s\alpha_p)(p^s t - p^s\beta_p) \in A_s$.
    The first estimate shows that $f_i = \lambda F_i(t)/C_i$ satisfies $f_i - g \in p^i A_s$, hence $f_i^a - g^a \in p^i A_s$.
    On the other hand, clearing the denominator of the word gives
    \[
        f_i^a = w\,\lambda^a C_i^{1 - a}\,\frac{H_{i, w}(t)}{C_i}
        \in p^i A_s w .
    \]
    Indeed $a \ge 1$, $1/C_i \in \mathbb{Z}_p$, and the second estimate in \eqref{eq:local-coefficient-estimate} applies to its last factor.
    Subtraction gives \eqref{eq:transfer-module} at this single finite stage.

    \emph{Passage to all primes.}
    At each prime assigned after stage $i$, the coefficients of $F_i$ are integral and its constant term is a unit.
    Thus $F_i$ takes a unit value at the primary argument there.
    Our residue exclusions give a unit at its guard as well.
    Consequently the prescribed valuation vectors for $U_i$ hold even at primes whose arguments had yet to be chosen.
    For a fixed nonzero polynomial $V \in \mathbb{Z}[t]$, the guard representatives $w_p$ tend to infinity and are $O(\log p)$ apart from the finite initial batch.
    It follows that $0 < |V(w_p)| < p$ for all sufficiently large $p$.
    Thus $V(\beta_p)$ is a unit at these primes.
    Applying this to numerators and denominators proves the guard assertion for every rational function.
    At primaries write a fixed such function as $t^k V(t)/W(t)$ with $V, W \in \mathbb{Z}[t]$ and $V(0)W(0) \ne 0$.
    At every sufficiently late prime these two constant terms are units and $\nu_p(\alpha_p) = 1$, giving primary valuation $k$.

    Finally, the profiles and the recursion give $a_n$ primary values zero at indices $j < 2n$, two at $j = 2n$, and one at $j > 2n$.
    Its guard valuations are zero and $a_n \in B$, so $a_n \in D$.
    Fix $m \ge 1$.
    For each $p \mid m$, write $p = p_j$.
    At every sufficiently late $q_n$ stage the prime has been assigned and $j < 2n$.
    Formula~\eqref{eq:atoms-profile} gives primary value zero for $a_n$ there, and \eqref{eq:profiles} gives zero for $b_{n + 1}$ and $q_n$ there.
    All three are also units at the guard, so both branches receive the refined target.
    Taking $n \ge \nu_p(m)$ gives integral values of $(a_{n + 1} - 1)/m$ at both branches.
    At all other primes $m$ is a unit, and the same quotient belongs to $B$ since $m$ is a nonzero rational constant.
    The intersection defining $D$ gives $(a_{n + 1} - 1)/m \in D$ as asserted.
\end{proof}
